% Part: proof-theory
% Chapter: natural-deduction
% Section: introduction

\documentclass[../../../include/open-logic-section]{subfiles}

\begin{document}

\olfileid{pt}{nat}{int}

\olsection{પરિચય}

સ્વાભાવિક નિગમન એવાં !!{proof}
તંત્રોનો પરિવાર છે જેમાં દરેક
તાર્કિક સંયોજક અથવા પરિમાણક
માટે અનુમાનોની એક જોડ હોય છે:
એક અનુમાન એ સંક્રિયકનો
``પરિચય'' કરાવે છે
(!!a{introduction} નિયમ)
અને બીજું તેનો
``નિવારણ'' કરે છે
(!!a{elimination} નિયમ).
દાખલા તરીકે, \Intro{\lif}માં
$!A \lif !B$ સ્વરૂપનાં
!!{formula}s નિષ્કર્ષરૂપે હોય છે,
જ્યારે~\Elim{\lif}માં
$!A \lif !B$ સ્વરૂપનાં
!!{formula}s
આધારવિધાનરૂપે હોય છે.

સ્વાભાવિક નિગમનનાં પ્રથમ બે
રૂપો ૧૯૩૦ના દાયકામાં ગેરહાર્ડ
ગેન્ટ્ઝન અને સ્તાનિસ્વાવ
યાશ્કોવ્સ્કીએ રજૂ કર્યાં હતાં.
સાબિતી-સિદ્ધાંતમાં ગેન્ટ્ઝનનું
તંત્ર વધુ ચર્ચાય છે, જ્યારે
શિક્ષણમાં સૌથી વધુ વપરાતાં
તંત્રો યાશ્કોવ્સ્કીની રીતમાંથી
વિકસ્યાં છે. બંનેમાં
\emph{ધારણાઓ}થી શરૂઆત
કરી શકાય છે: એવાં
!!{formula}s જેમને
!!{prove}d કરવાં જરૂરી
નથી, પણ જે બીજા
!!{formula}sને !!{proving}
કરવા વાપરી શકાય છે.
આખી !!{proof} આવી
ધારણાઓ પર આધારિત હોઈ
શકે છે. તે ધારણાઓ
અસાબિત હોવાથી !!{proof}
પોતાના નિષ્કર્ષને
(પોતાના \emph{અંતિમ-!!{formula}}ને)
સ્થાપિત કરતી નથી;
માત્ર એટલું સ્થાપિત
કરે છે કે નિષ્કર્ષ
ધારણાઓમાંથી ફલિત થાય છે.

કેટલાક અનુમાન-નિયમોના
નિષ્કર્ષો પછી કેટલીક ધારણાઓ
પર આધારિત રહેતા નથી:
નિયમો એ ધારણાઓને
``મુક્ત'' કરે છે.
દાખલા તરીકે,
\Intro{\lif} નિયમ
ધારણા~$!A$માંથી
નિષ્કર્ષ~$!B$ની
!!a{proof} લે છે અને
$!A \lif !B$ નિષ્કર્ષ આપે છે.
$!A \lif !B$માં પૂરી થતી
!!{proof} હવે~$!A$ પર
આધારિત નથી; ધારણા~$!A$
મુક્ત થાય છે.
ગેન્ટ્ઝનના તંત્રોમાં
!!{proof}s એ
!!{formula}sનાં વૃક્ષો છે;
ધારણાઓ સૌથી ઉપરનાં
!!{formula}s છે;
અને~\Intro{\lif} જેવા
નિયમો પરનાં ચિહ્નો કઈ
ધારણાઓ મુક્ત થાય છે
તે બતાવે છે.
યાશ્કોવ્સ્કીના તંત્રોમાં
ધારણા પર આધારિત
!!{proof}નો ભાગ અલગ
દેખાડાય છે, દાખલા તરીકે,
ઊભી રેખાની પાછળ
ગોઠવીને અથવા ચોકઠામાં
મૂકીને. એ ચોકઠાની
પહેલી પંક્તિમાં મુક્ત
થતી ધારણા~$!A$ અને
છેલ્લી પંક્તિમાં
અનુસરણનું પરિણામ~$!B$
હોય છે; \Intro{\lif}નો
નિષ્કર્ષ~$!A \lif !B$
ચોકઠાની બહાર હોય છે.

આ તંત્રો ``સ્વાભાવિક''
છે, કારણ કે તેમનાં
અનુમાન-નિયમો આપણે
(ઓછામાં ઓછા ગણિતશાસ્ત્રીઓ)
સ્વાભાવિક રીતે કેવી રીતે
તર્ક કરીએ છીએ તે દર્શાવે છે.
શરતી પરિણામ સ્થાપવા માટે
આપણે તેનું પૂર્વપદ
ધારીને પરિણામપદ સાબિત
કરીએ છીએ: એ~\Intro{\lif}.
$!A \lif !B$ અને તેનું
પૂર્વપદ~$!A$ જાણીએ
ત્યારે~$!B$ તારવીએ છીએ:
એ~\Elim{\lif}.
વિકલ્પ~$!A \lor !B$
ધારતાં કોઈ નિષ્કર્ષ
મળે છે તે કિસ્સાવાર
સાબિતીથી બતાવીએ છીએ:
એ~\Elim{\lor}.
સામાન્ય દાવો
$\lforall[x][!A(x)]$
સાબિત કરવા કોઈ મનસ્વી~$c$
માટે~$!A(c)$ બતાવીએ છીએ:
એ~\Intro{\lforall}.
અને આવી રીતે આગળ.

દાખલા તરીકે, ગેન્ટ્ઝનની
શૈલીના સ્વાભાવિક નિગમનમાં
$!A \lif (!A \lor !B)$ની
આ સરળ !!{proof} જુઓ:
\[
\AxiomC{$\Discharge{!A}{x}$}
\RightLabel{\Intro{\lor}}
\UnaryInfC{$!A \lor !B$}
\DischargeRule{\Intro{\lif}}{x}
\UnaryInfC{$!A \lif (!A \lor !B)$}
\DisplayProof
\]
તે ધારણા~$!A$થી શરૂ થાય છે,
\Intro{\lor}થી~$!A
\lor !B$ તારવે છે,
અને પછી~\Intro{\lif}થી
$!A \lif (!A \lor !B)$
તારવે છે. \Intro{\lif}
અનુમાન ધારણા~$!A$ને
મુક્ત કરે છે; ચિહ્ન~$x$
એ દર્શાવે છે.

સાબિતી-સિદ્ધાંતકારો
!!{formula}sનાં વૃક્ષો
વાપરતાં ગેન્ટ્ઝનનાં મૂળ
તંત્રોને પસંદ કરે છે,
છતાં એક અન્ય રૂપ પણ
સાબિતી-સિદ્ધાંતમાં
મહત્ત્વનું છે.
ધારણાઓ~$\Gamma$માંથી
$!A$ની !!^a{proof}
બતાવે છે કે~$!A$ એ
$\Gamma$માંથી ફલિત થાય છે,
અર્થાત્~$\Gamma \Entails !A$.
સ્વાભાવિક નિગમનના
અનુમાન-નિયમોને આ પ્રકારની
ફલિતતાઓ વિશેના નિયમો
તરીકે પણ વાંચી શકાય છે.
દાખલા તરીકે,
\Intro{\lif} કહે છે કે
જો $\Gamma + !A
\Entails !B$, તો
$\Gamma \Entails !A \lif !B$.
\sourcecorrection{OLMD-227}{મૂળમાં અનુસરણ-પરિચયનું નિષ્કર્ષ $\Gamma \Entails !B$ લખાયું છે; ધારણા $!A$ મુક્ત કર્યા પછી યોગ્ય નિષ્કર્ષ $\Gamma \Entails !A \lif !B$ છે.}
આથી નિયમોને સીધા
એ રીતે ઘડવું સ્વાભાવિક છે
કે તેઓ આધારવિધાન અને
નિષ્કર્ષ બંનેમાં
ધારણાઓ તથા તેમનામાંથી
મળતા પરિણામ પર કામ કરે;
અર્થાત્ તેઓ
$\Gamma \Sequent !A$
સિક્વન્ટો પર કામ કરે.
ઉપરની સરળ સાબિતી
આવા તંત્રમાં આ રીતે
દેખાય:
\[
\Axiom$x:!A \fCenter !A$
\RightLabel{\Intro{\lor}}
\UnaryInf$x:!A \fCenter !A \lor !B$
\DischargeRule{\Intro{\lif}}{x}
\UnaryInf$\fCenter !A \lif (!A \lor !B)$
\DisplayProof
\]
અહીં નિયમો એ જ છે:
તેઓ~$\Sequent$ની જમણી
બાજુના !!{formula} પર કામ કરે છે;
ડાબી બાજુનાં !!{formula}s
દરેક પગલે આધારરૂપ
ધારણાઓની યાદી આપે છે.

!!{introduction}/!!{elimination}
નિયમોની બધી જોડો
એકલી મળીને શાસ્ત્રીય
તર્ક માટે પૂર્ણ નથી.
$\lfalse$ને લગતા
બે વધારાના નિયમો
ઉમેરવા પડે છે.
પહેલો ``અંતઃપ્રજ્ઞાવાદી
અસંગતિ-નિયમ''~\FalseInt,
અથવા ``વિસ્ફોટ'' છે:
$\lfalse$માંથી
કંઈ પણ તારી શકાય.
બીજો પરોક્ષ સાબિતીનો
શાસ્ત્રીય સિદ્ધાંત~\FalseCl{}
(અથવા ``શાસ્ત્રીય
અસંગતિ-નિયમ'') છે:
જો~$\lnot !A$માંથી
$\lfalse$ !!{prove}
કરી શકાય, તો~$!A$
!!{prove}d થયું છે.
જો~\FalseCl કાઢી નાખીએ,
તો અંતઃપ્રજ્ઞાવાદી
તર્ક માટે યોગ્ય તંત્ર મળે છે.
બંને કાઢી નાખીએ,
તો ``લઘુતમ તર્ક''
નામે ઓળખાતું તંત્ર મળે છે.

હવે શાસ્ત્રીય અને
અંતઃપ્રજ્ઞાવાદી તર્ક માટે
ગેન્ટ્ઝનની શૈલીનું
સ્વાભાવિક નિગમન જોઈશું.
આ તંત્રો~\Log{N1c}
અને~\Log{N1i} છે.
એકલાં !!{formula}s~$!A$
ને બદલે સિક્વન્ટો
$\Gamma \Sequent !A$
પર કામ કરતાં અનુરૂપ
તંત્રો~\Log{N2c}
અને~\Log{N2i} છે.

\end{document}
